\section{The Check You Turned Off}
\label{sec:ch10-the-check-you-turned-off}
\index{satisfiability!what the router does with it off|(}%

Chapter~\ref{ch:composition} ended a section by refusing to guess. It had
composed a deliberately unsatisfiable graph with
\code{--disable-resolvability-validation}, found that the config it produced
carries a client schema \code{===} identical to a healthy one, and said that
what a router does with such a file was not something it was going to describe
without watching. One request answers it. Here is the request.

The setup is chapter~\ref{ch:composition}'s \code{unsatisfiable-key} case
unchanged. Mosaic's \code{Product} keeps its key and gains one argument,
\code{resolvable: false}, and the composer answers with four unresolvable-path
errors and no output. With the flag it produces a config instead, in which the
only thing recording that anything is wrong is one property on one key:

\begin{minted}{json}
{ "typeName": "Product", "selectionSet": "id", "disableEntityResolver": true }
\end{minted}

I gave that file to a router. Five things happen, in this order.

\subsection*{It starts, and says nothing}

\index{router!silent on a broken graph}%
The router comes up. \code{/health} answers 200. The startup log is eleven
lines again, and it differs from
section~\ref{sec:ch10-one-process-two-files}'s in exactly one of them: this
config does not turn watching on, so where Mosaic's router says it is watching
the execution config, this one says that a static config was provided and
updating it needs a restart. The other ten are identical, warning about the
missing graph token and about development mode, and carrying nothing whatsoever
about the graph it just loaded. I counted the warning lines mentioning
resolvability, unreachable fields, entities or subgraphs, and there are none.

That is the first surprise and the largest. The router has, in one file, a
schema promising four fields on \code{Product} and a routing table saying the
only service holding them may not be entered. Both halves are in front of it at
startup. It could compare them. It does not.

\subsection*{Most of the graph works}

Ask for something that stays inside Catalog and it is fine:

\begin{minted}{graphql}
{ browseProducts(first: 2) { nodes { title sku } } }
\end{minted}

Two products, HTTP 200, no errors. Which makes this worse rather than better.
A deployment of this config passes a smoke test that reads a list, and passes
whatever health check you have, and serves the majority of a storefront
correctly. What it cannot do is the thing federation was for.

\subsection*{Crossing the boundary is a 500}

\begin{minted}{graphql}
{ browseProducts(first: 2) { nodes { title price { amount } } } }
\end{minted}

\begin{minted}[fontsize=\footnotesize]{text}
HTTP 500

{"errors":[{"message":"internal server error"}]}
\end{minted}

\index{errors!internal server error@\code{internal server error}}%
No field name. No subgraph name. No hint that the graph was composed with a
check switched off. A client developer receiving this has no route from the
message to the cause, and the operator reading it has one only if they already
suspect the answer.

Note also that this one is a 500 where
section~\ref{sec:ch10-sixty-six-settings}'s unreachable subgraph was a 200. The
distinction is real: an upstream that will not answer is a data problem the
GraphQL response can carry, and a query the router cannot plan is not a query at
all.

\subsection*{The log knows, and names the wrong type}

The router's own log has the cause:

\begin{minted}[fontsize=\footnotesize,breaklines]{text}
1 error occurred: * printOperation planner id: 1: validation failed: external: Cannot query field "price" on type "Query"., locations: [], path: [query]
\end{minted}

\index{errors!that name the wrong type}%
\code{price} is not a field of \code{Query}. It never has been, in any schema in
this book, and nothing the client sent suggests it should be. Read the message
twice before dismissing it, because it is not nonsense: it is about a document
the planner wrote, not the one the client sent. Denied the only route into
Mosaic, the planner built a subgraph request that asks for \code{price} at the
root, and that document failed validation against Mosaic's schema. The words
\code{printOperation planner} at the front are the tell that the subject is the
router's own output.

Figure~\ref{fig:ch10-where-the-failure-went} is the whole trade the flag makes.

\begin{figure}[htbp]
  \centering
  % Where --disable-resolvability-validation moves the failure. Referenced from
% chapter 10. Bare tikzpicture; the figure environment, caption and label live
% at the call site.
%
% Two columns for the same fault, and one horizontal rule. Everything above the
% rule happens on a laptop with both schemas open; everything below it happens
% in production. The columns run top to bottom rather than left to right
% because the right-hand one has three outcomes and no room for them sideways.
\begin{tikzpicture}[
    font=\small,
    stage/.style={draw, rounded corners=2pt, minimum width=46mm,
                  minimum height=8mm, font=\scriptsize, align=center},
    result/.style={draw, rounded corners=2pt, minimum width=52mm,
                   minimum height=12mm, font=\scriptsize, align=center,
                   fill=black!8},
    quiet/.style={draw, rounded corners=2pt, minimum width=52mm,
                  minimum height=11mm, font=\scriptsize, align=center,
                  dashed},
    flow/.style={-{Stealth[length=2mm]}},
    lbl/.style={font=\scriptsize, align=center, inner sep=1pt},
    divider/.style={gray, dotted, thick}
  ]

  % -- the two columns -------------------------------------------------------
  \node[lbl] at (3.5,0.5) {\textbf{Composed as it stands}};
  \node[lbl] at (10.5,0.5) {\textbf{Composed with the check off}};

  \node[stage] (cmdA) at (3.5,-0.4)  {\code{wgc router compose}};
  \node[stage] (cmdB) at (10.5,-0.4)
    {\code{wgc router compose}\\\code{--disable-resolvability-validation}};

  \node[result] (errA) at (3.5,-2.0)
    {4 errors, exit 1\\each names a field\\and the path a client takes to it};
  \node[result] (cfgB) at (10.5,-2.0)
    {a config, exit 0\\one property differs\\from a healthy one};

  \draw[flow] (cmdA) -- (errA);
  \draw[flow] (cmdB) -- (cfgB);

  % -- the divide ------------------------------------------------------------
  \draw[divider] (0.3,-3.1) -- (13.7,-3.1);
  \node[lbl, gray, anchor=south east] at (13.7,-3.03) {build time};
  \node[lbl, gray, anchor=north east] at (13.7,-3.17) {request time};

  % -- what each one leaves you with -----------------------------------------
  \node[lbl, align=center] (noneA) at (3.5,-4.1)
    {nothing is deployed,\\because nothing\\was produced};

  \node[quiet] (routerB) at (10.5,-4.1)
    {the router starts\\no warning about the graph};
  \draw[flow] (cfgB) -- (routerB);

  \node[result] (okB) at (10.5,-5.8)
    {a query inside one subgraph\\HTTP 200};
  \node[result] (badB) at (10.5,-7.4)
    {a query across the boundary\\HTTP 500,
     \code{internal server error}};

  % Two outcomes of the same router, not a sequence, so the second one is
  % routed around the first rather than through it.
  \draw[flow] (routerB.south) -- (okB.north);
  \draw[flow] (routerB.west) -- (7.4,-4.1) -- (7.4,-7.4) -- (badB.west);

\end{tikzpicture}

  \caption{The same fault, twice. Without the flag it is a build-time message
  that names the field, the path a client would take to it and why the walk
  failed. With it, everything to the left of the router is quiet, and what
  reaches the operator instead is a 500 with no field in it and a log line about
  a document nobody wrote.}
  \label{fig:ch10-where-the-failure-went}
\end{figure}

\subsection*{It fails before anything is called}

\index{query plan!failing at plan time}%
The last measurement is the one that surprised me most, and it is a timing. The
request logger records the failed request at \code{"latency": 0.000822445},
which is eight hundred microseconds. No network call fits in that.

The confirmation is stronger than the timing: the same 500 comes back with
\code{X-WG-Skip-Loader: true} set, and that header means the router does not
execute the plan at all. So the failure is not a subgraph refusing the request.
It is the router failing to build one.

Which is, in a bleak way, the good news buried in this section. The graph is
broken deterministically rather than intermittently, every query that crosses
the boundary fails the same way every time, and it fails before it costs a
subgraph anything. There is no partially-served response and no half-written
data. What there is instead is a class of query that has stopped working with no
message anybody can act on.

\subsection*{Both halves are asserted}

\index{verification!router cases}%
Chapter~\ref{ch:composition} made its composition errors a gate, on the grounds
that a printed error message goes stale as easily as a schema. Everything in
this section is the same kind of claim, so it is the same kind of gate:
\code{scripts/router-cases.mjs}, called by both verification scripts, one node
implementation for the reason chapter~\ref{ch:composition} gives.

\begin{minted}[fontsize=\footnotesize]{text}
node scripts/router-cases.mjs --list

config-beats-env    a value in config.yaml wins over the environment variable for it
localhost-fallback  localhost in a routing URL reaches the host, until you turn that off
resolvability-off   a graph composed with --disable-resolvability-validation, in front of a router
\end{minted}

Each case starts real containers against a real config and each carries its own
control, which is the part I would keep if I kept one thing. The precedence case
fails if \code{LOG\_LEVEL=debug} ever stops producing debug lines on its own;
the fallback case fails if the query succeeds with the fallback off. Without
those, both would go on passing against a router that had simply stopped
answering, and a green gate that cannot fail is worse than no gate.

\index{verification!a counting trap}%
The precedence case earned its control immediately. My first version counted
zero debug lines in both runs and would have reported the finding backwards. The
router colours its log, so the bytes around the word are escape sequences rather
than the spaces I was matching on, and every level check returned zero. The
control caught it in one run. Chapter~\ref{ch:composition} spent a decision on
this exact lesson and I still walked into it: a count is of a thing, not of what
your pattern happened to match.

\medskip
So the flag does what its help text warns, and now the warning can be stated in
full. It converts a build-time failure that names four fields, the query path to
each of them and the key that was missing, into a runtime failure that names
nothing, arrives after deployment, and reaches whoever is on call rather than
whoever caused it. Chapter~\ref{ch:composition} called that the worst available
outcome without having seen it. It was right.

That is the router as a process: what it loads, what it costs, and what it does
with a graph you told it not to check.
Chapter~\ref{ch:entity-resolution-done-right} goes back to the subgraphs and
makes the second hop do more work than a key lookup.

\index{satisfiability!what the router does with it off|)}%
